\chapter{Cointegration and VECM}
\label{cap:coint}
% ── index ──────────────────────────────────────────────────────
\index{cointegration|textbf}
\index{VECM|textbf}
\index{error correction model|see{VECM}}
\index{Johansen, test|textbf}
\index{johansen@\texttt{johansen()}|textbf}
\index{vecm@\texttt{vecm()}|textbf}
\index{long-run relationship}

% ─────────────────────────────────────────────────────────────────────────────
\section{Long-run Relationship Between I(1) Series}

Two $I(1)$ series are \textbf{cointegrated} if there exists a linear combination
between them that is $I(0)$. Formally, $y_t$ and $x_t$ are cointegrated if
there exists $\beta$ such that:
\[
  u_t = y_t - \beta x_t \sim I(0)
\]

The intuition is that of a \emph{leash}: series that individually follow
random walks may be attracted by a long-run equilibrium force. Asset prices
in the same market, exchange rates and purchasing power parity, short- and
long-term interest rates are classic examples.

Non-cointegrated $I(1)$ series diverge indefinitely — level regression produces
\textbf{spurious regression}: high $R^2$ and significant $t$-statistics with
no real causal relationship.

% ─────────────────────────────────────────────────────────────────────────────
\section{Engle-Granger Test}

The Engle-Granger (1987) procedure tests cointegration in two steps:

\begin{enumerate}
  \item Estimate the \textbf{long-run relationship} by OLS:
        $\hat u_t = y_t - \hat\beta x_t$
  \item Apply the \textbf{ADF test} to the residuals $\hat u_t$ — under $H_0$
        the residuals have a unit root (no cointegration).
\end{enumerate}

The critical values of the test are more negative than those of the standard ADF,
since $\hat u_t$ is an estimated residual, not a directly observed series.

% ─────────────────────────────────────────────────────────────────────────────
\section{Error Correction Model (VECM)}

If $y_t$ and $x_t$ are cointegrated, the VECM (\emph{Vector Error Correction
Model}) decomposes the dynamics into a short-run component and adjustment to
the long-run equilibrium:

\[
  \Delta \mathbf{y}_t
  = \boldsymbol{\alpha} \underbrace{(\boldsymbol{\beta}' \mathbf{y}_{t-1})}_{\text{deviation from equilibrium}}
  + \sum_{j=1}^{p-1} \mathbf{\Gamma}_j \Delta \mathbf{y}_{t-j}
  + \boldsymbol{\varepsilon}_t
\]

\begin{itemize}
  \item $\boldsymbol{\beta}$ — \textbf{cointegrating vector}: defines the long-run
        equilibrium relationship $\boldsymbol{\beta}'\mathbf{y}_t = 0$.
  \item $\boldsymbol{\alpha}$ — \textbf{adjustment coefficients}: speed at which
        each variable corrects deviations from equilibrium. A negative sign in
        $\alpha_y$ means $y$ falls when it is above equilibrium.
  \item $\mathbf{\Gamma}_j$ — short-run dynamics.
\end{itemize}

The VECM is estimated by ML via the Johansen procedure, which tests the number
of cointegrating vectors (\emph{rank}) via the trace statistic or maximum eigenvalue.

% ─────────────────────────────────────────────────────────────────────────────
\section{DGP Verification}

\subsection*{Case 1 — cointegrated: $y_t = 2 x_t + \varepsilon_t$}

$x_t$ is a random walk ($I(1)$); $\varepsilon_t \sim I(0)$. Therefore
$y_t - 2 x_t = \varepsilon_t \sim I(0)$: the series are cointegrated with
vector $[1,\; -2]$.

\begin{lstlisting}[caption={Cointegrated DGP and Engle-Granger test}]
seed(42)
input df
id
1
end
for i in 1..=9 { df = append(df, df) }
df |> filter(_n <= 300)
generate df e1 = rnormal()
generate df e2 = rnormal()
generate df x  = cumsum(e1)
generate df y  = 2.0*x + e2
coint(df, y, x)
\end{lstlisting}

\begin{lstlisting}[language={}, basicstyle=\ttfamily\small,
  frame=none, numbers=none, backgroundcolor=\color{codbg},
  xleftmargin=0pt, xrightmargin=0pt, breaklines=false]
============= Engle-Granger Cointegration Test =============
  Variables: y, x
  H₀: no cointegration
  ADF statistic:    -6.4808
  Critical values:  1%=-3.900  5%=-3.340  10%=-3.040
  Cointegrating vector: [0.4778, 2.0002]
  Conclusion: REJECT H₀ — cointegrated series
\end{lstlisting}

The statistic $-6{,}48$ is well below the 1\% critical value
($-3{,}90$). The estimated cointegrating vector $[0{,}48;\; 2{,}00]$ recovers
$\beta = 2{,}0$ with precision.

\subsection*{Case 2 — not cointegrated: two independent random walks}

\begin{lstlisting}[caption={Two random walks — no cointegration}]
generate df y = cumsum(e2)
coint(df, y, x)
\end{lstlisting}

\begin{lstlisting}[language={}, basicstyle=\ttfamily\small,
  frame=none, numbers=none, backgroundcolor=\color{codbg},
  xleftmargin=0pt, xrightmargin=0pt, breaklines=false]
============= Engle-Granger Cointegration Test =============
  Variables: y, x
  H₀: no cointegration
  ADF statistic:    -3.1058
  Critical values:  1%=-3.900  5%=-3.340  10%=-3.040
  Conclusion: Fail to reject H₀ — no cointegration
\end{lstlisting}

\subsection*{VECM — adjustment dynamics}

With cointegration confirmed, we estimate the VECM:

\begin{lstlisting}[caption={VECM with rank 1}]
generate df y = 2.0*x + e2
let m = vecm(df, y, x, lags=1)
display m
\end{lstlisting}

\begin{lstlisting}[language={}, basicstyle=\ttfamily\small,
  frame=none, numbers=none, backgroundcolor=\color{codbg},
  xleftmargin=0pt, xrightmargin=0pt, breaklines=false]
======================== VECM (Johansen ML) - Rank 1 =========================
No. Variables:                2
Observations:               299

------------------------ Cointegration Vectors (Beta) ------------------------
[ 3.6175, -7.2357 ]  →  normalized: [ 1, -2.000 ]

---------------------- Adjustment Coefficients (Alpha) -----------------------
  y:  -0.3133   (corrects deviations from equilibrium)
  x:  -0.0179   (nearly exogenous — minimal adjustment)

----------------------- Johansen Eigenvalues (Lambda) ------------------------
    λ₁ = 0.5017     λ₂ = 0.0004
==============================================================================
\end{lstlisting}

The normalized cointegrating vector $[1,\; -2{,}00]$ recovers the true
$\beta = 2{,}0$. The adjustment coefficient $\hat\alpha_y = -0{,}313$: when
$y_{t-1} > 2 x_{t-1}$ (above equilibrium), $\Delta y_t$ is negative —
$y$ corrects downward. $x$ barely adjusts ($\hat\alpha_x \approx -0{,}02$),
since it is the exogenous variable that defines the equilibrium.

% ─────────────────────────────────────────────────────────────────────────────
\section{Syntax}

\begin{lstlisting}
coint(df, y, x)
coint(df, y, x, lags=4)
vecm(df, y, x, lags=1)
vecm(df, y1, y2, y3, lags=2)
\end{lstlisting}

\begin{lstlisting}[caption={Typical workflow — exchange rate and parity}]
load "cambio.csv" as df

adf(df, spot)
adf(df, forward)

coint(df, spot, forward)

// if cointegrated:
let m = vecm(df, spot, forward, lags=2)
display m
\end{lstlisting}

\begin{nota}
  VECM assumes all variables are $I(1)$ and cointegrated.
  Verify with ADF before estimating. If the series are $I(1)$ but not
  cointegrated, use the VAR in first differences (ch.~\ref{cap:var}).
\end{nota}
